\documentclass[../main.tex]{subfiles}
\begin{document}
%==========================================TIÊU ĐỀ TẬP========================================
\begin{table}[h]
    \begin{adjustwidth}{-1cm}{}
        \begin{tabular}{>{\centering\arraybackslash}p{6cm}|>{\raggedright\arraybackslash}p{12.5cm}}
            \multirow{5}{6cm}
            % Cột bên trái: ÉPISODE và số tập
            {\\[-40pt]
                \centering
                {\fontsize{20pt}{20pt}\selectfont{\outlinetext[1]{eptype}{white}{\flaregothic ÉPISODE}}\\[12pt]
                    {\fontsize{72pt}{72pt}\selectfont{\outlinetext[1]{epnum}{white}{\digi C}}}\\[6pt]       % use for normal/special episodes
                    % {\fontsize{36pt}{36pt}\selectfont{\outlinetext[1]{epnum}{white}{\burbank SÁU}}}\\[6pt]     % use for custom episodes
                    {\fontsize{20pt}{20pt}\selectfont{\outlinetext[1]{epnum}{white}{\cafeteria (KL-C)}}}\color{black}}
                \\[18pt]
                % {\fontsize{14pt}{14pt}\selectfont{\outlinetext[1]{epattr}{white}{\flaregothic [I]}}}\\
                % {\fontsize{18pt}{18pt}\selectfont{\outlinetext[1]{epattr}{white}{\flaregothic [II]}}}\\
                \color{black}
            }
             &
            % Dòng thứ nhất: tên tập tiếng Nhật
            {\fotskip\fontsize{18pt}{18pt}\selectfont \ruby{脱}{だつ}\ruby{走}{そう}}                                                  \\[6pt]
            % Dòng thứ hai: tên tập tiếng Anh
             & {\cafeteria\fontsize{18pt}{18pt}\selectfont Escape from holoX}                                                    \\[4pt]
            % Dòng thứ ba: tên tập tiếng Việt
             & {\vnmsans\fontsize{18pt}{18pt}\selectfont Tẩu thoát khỏi holoX}                                                   \\[8pt]
            % Dòng thứ tư: Nhân vật xuất hiện
             & {\caviar{\fontsize{14pt}{14pt}\selectfont Nhân vật: \emp{koyori} \emp{laplus} \emp{lui} \emp{chloe} \emp{iroha}}} \\[8pt]
            % Dòng cuối cùng: Thời gian diễn ra của tập (thường sẽ là ngày phát hành)
            %  & {\sen\fontsize{16pt}{16pt}\selectfont Mốc thời gian: DD/MM/YYYY}  \\[18pt]
        \end{tabular}\\[8pt]
    \end{adjustwidth}
\end{table}
%=========================================TRUYỆN CHÍNH========================================
\color{black}

\txtbx[2]{
    {\color{vol} \uline{Tóm tắt:}} Hơn một tháng trôi qua kể từ đêm tái sinh định mệnh, ba mẹ con Mi, Tô và Ki trải qua chuỗi ngày huấn luyện khắc nghiệt dưới trướng năm thành viên của tổ chức bí mật holoX. Kết quả kiểm tra năng lực cho thấy các chỉ số sinh học phi thường: Mi sở hữu trí tuệ và sức mạnh gấp mười lần cả gia đình Hikawa cộng lại; Tô nhân đôi thể lực dũng mãnh nhưng trí lực giảm phân nửa; còn Ki nhân đôi tư duy logic siêu việt dù sức mạnh suy giảm. Bằng sự nhẫn nại và lòng trung thành nguyên bản, ba mẹ con lần lượt vượt qua các bài kiểm tra của Koyori, Iroha, Lui, Chloe và La+. Tuy nhiên, sự phục tùng bên ngoài chỉ là lớp vỏ bọc hoàn hảo nhằm che giấu một kế hoạch táo bạo. Khi bài kiểm tra tối hậu — trận giả chiến thuật tại bãi cát ven biển sa mạc — được phát động, thời cơ ngàn năm có một đã mở ra để ba sinh mệnh tìm đường quay về mái ấm thực sự tại Kawachi.
}

\begin{center}
    \uline{Phòng đo đạc thể chất ngầm của \textit{holoX}.}
\end{center}

Sau khi được cấp phát ba bộ trang phục huấn luyện đặc chế bằng sợi vải co giãn màu xám tro - có khoét hai lỗ nhỏ trên đỉnh mũ trùm cho đôi tai cún và một rãnh hở phía sau thắt lưng cho chiếc đuôi xù - ba mẹ con Mi, Tô và Ki bắt đầu bước vào đợt kiểm tra năng lực toàn diện đầu tiên.

Koyori đứng trước màn hình hiển thị không gian ba chiều, đeo cặp kính bảo hộ và cầm bút quang học ghi chép liên tục. Đôi tai sói trên đỉnh đầu cô nàng giật giật nhạy bén theo từng đợt sóng xung điện phát ra từ máy quét sinh trắc học: ``Nào, bắt đầu từ người mẹ trước nhé. Mi, xin hãy dùng toàn lực đấm thẳng vào tấm bia hấp thụ lực bằng hợp kim titan phía trước.''

Người phụ nữ mang vóc dáng thiếu phụ đằm thắm khẽ gật đầu. Mi bước lên một bước, hít sâu một hơi. Đôi tai chó trên đỉnh đầu khẽ áp sát vào làn tóc tro bạc, rồi cánh tay thon thả tung ra một cú đấm thẳng thớm, chuẩn xác vào tâm tấm bia.

*RẦM*

Cả sàn phòng thí nghiệm rung chuyển nhẹ. Đồng hồ đo áp lực nhảy số điên cuồng, dừng lại ở con số mười triệu Newton. Không chỉ có vậy, bài kiểm tra phản xạ thần kinh và tính toán ma trận ma pháp tiếp theo của Mi cũng đạt điểm số tuyệt đối trong sự ngỡ ngàng của mọi người.

Koyori đẩy gọng kính, nuốt nước bọt:

``Kufufu\dots\ Không thể tin nổi. Sức mạnh cơ bắp và chỉ số tư duy của Mi đạt mức cân bằng hoàn mỹ. Nếu lấy chỉ số trung bình của toàn bộ sáu thành viên gia đình Hikawa cộng lại làm hệ quy chiếu chuẩn, thì năng lực tổng hợp của Mi gấp đúng mười lần.''

``Mười lần cả nhà cộng lại sao?'' La+ ngồi trên vai chiếc ghế xoay, hai mắt tròn xoe kinh ngạc. ``Thế còn hai đứa con gái thì sao hả Koyori?''

Koyori vẫy tay gọi cô con gái lớn: ``Đến lượt Tô. Hãy thử nhấc khối tạ áp suất kia lên xem nào.''

Tô mở to đôi mắt đen láy tràn đầy vẻ tinh nghịch. Chiếc đuôi dài cong vút như lưỡi câu phía sau quẫy tít mù gõ cồm cộp xuống sàn nỉ. Cô nàng bước tới, chẳng cần dùng đến toàn bộ cánh tay, cô chỉ cần dùng\dots\ đúng một ngón tay săn chắc đã nhấc bổng khối tạ nặng hai tấn lên không trung nhẹ như một món đồ chơi.

Nhưng khi Koyori chuyển sang bài kiểm tra tính toán hình học không gian, Tô liền ngồi vò đầu bứt tai. Cô nàng cắn nát ngọn bút chì, mặt nhăn nhó như sắp khóc, thỉnh thoảng lại đưa tay lên gãi gãi đôi tai chó vểnh cao trên đầu: ``Ư\dots\ Khó quá\dots\ Nhìn mấy cái hình này nhức đầu quá mẹ ơi\dots''

Mi khẽ mỉm cười, đưa tay xoa đầu dỗ dành con gái.

Koyori nhìn vào biểu đồ phân tích, gật gù ghi chú: ``Thực thể Tô thừa hưởng trọn vẹn gen chó cỏ săn mồi dũng mãnh. Sức mạnh thể chất và tốc độ phản xạ của con bé gấp đôi người mẹ, tức là gấp hai mươi lần cả nhà Hikawa gộp lại. Thế nhưng, đổi lại thì năng lực tư duy trừu tượng bị giảm đi một nửa so với Mi, chỉ ở mức gấp năm lần người thường. Một chiến binh tiên phong thuần túy thiên về bản năng và hành động trực giác.''

Cuối cùng là cô con gái út Ki.

Cô bé mười lăm tuổi với vóc dáng mảnh dẻ bước lên bục kiểm tra. Khi đối mặt với tấm bia lực, cú đấm của Ki chỉ tạo ra một lực bằng một nửa của mẹ Mi, tương đương năm lần người thường.

Thế nhưng, khoảnh khắc Ki ngồi vào bàn điều khiển bài kiểm tra trí tuệ và giải mã mật mã lượng tử, cả căn phòng thí nghiệm bỗng chốc rơi vào sự tĩnh lặng tuyệt đối.

Những ngón tay nhỏ nhắn, trắng nõn nà của Ki lướt trên bàn phím ảo nhanh đến mức chỉ còn lại những tàn ảnh lấp loáng. Chưa đầy bốn mươi giây, cô bé đã bẻ khóa toàn bộ mười hai tầng tường lửa mô phỏng của căn cứ holoX, đồng thời chỉ ra ba lỗ hổng logic trong thuật toán ngưng tụ năng lượng của Koyori.

Đôi tai cún phủ lông tơ rủ nhẹ hai bên trán của Ki khẽ giật nhẹ, giọng nói trong veo cất lên đều đều: ``Phương trình số bảy ở góc màn hình có sai số không phẩy không ba phần trăm ở hệ số tỏa nhiệt. Nếu không hiệu chỉnh, lò phản ứng sẽ phát nổ sau ba trăm giờ vận hành.''

Nàng Sói hồng há hốc mồm, chiếc bút quang học tuột khỏi tay rơi đánh cộp xuống nền đất: ``T-Trời đất ơi\dots\ Chỉ số trí tuệ và năng lực tính toán của Ki gấp đôi Mi, tức là gấp hai mươi lần cả nhà Hikawa cộng lại. Con bé sở hữu bộ não của một siêu máy tính sinh học. Một bộ ba tương hỗ hoàn mỹ ngoài sức tưởng tượng: người mẹ là hạt nhân chỉ huy toàn năng, người chị là mũi nhọn càn quét vô song, và cô em út là bộ não chiến lược tối thượng.''

La+ Darknesss nghe xong thì nhảy cẫng lên reo hò đắc chí, chiếc quạt giấy vung vẩy phần phật: ``Bwahahaha! Tuyệt vời! Quá sức vĩ đại! Với bộ ba cận vệ siêu đẳng này, ngày mà La+ Darknesss thống trị thế giới sẽ không còn xa nữa! Bắt đầu huấn luyện thực chiến ngay lập tức!''

\ct

Suốt hơn một tháng sau đó, căn cứ ngầm của \textit{holoX} trở thành một võ đài huấn luyện không ngừng nghỉ.

Để biến ba sinh mệnh vốn quen với cuộc sống canh xưởng mộc mạc thành những chiến binh tinh nhuệ, năm thành viên của holoX lần lượt đứng ra đảm nhiệm từng học phần chuyên biệt.

Bài kiểm tra đầu tiên về trí tuệ và phản xạ thông tin do Koyori phụ trách. Ki hoàn thành xuất sắc mọi giáo trình trong sự thán phục của nàng Sói hồng, trong khi Tô liên tục làm nổ các ống nghiệm vì lỡ tay bóp quá mạnh, khiến phòng lab ngập trong khói hồng và những tiếng la hét dở khóc dở cười.

Bài kiểm tra thứ hai về thể thuật và kiếm đạo do Iroha đích thân chỉ dạy. Giữa võ đường lót sàn gỗ, Tô thể hiện bản năng cận chiến kinh hoàng. Cầm thanh kiếm gỗ trên tay, cô nàng vung một nhát chém tạo ra luồng kình phong mạnh đến mức bạt văng tấm bình phong ngăn cách, ép Iroha phải lùi lại năm bước để giữ thăng bằng.

Nàng samurai tóc vàng vừa xoa cổ tay tê dại vừa tấm tắc ngợi khen: ``Nội lực thâm hậu và bộ pháp dũng mãnh phi thường, de gozaru. Dù chưa từng học qua kiếm phổ, nhưng trực giác chiến trận của Tô-dono quả thật đáng kinh ngạc.''

Bài kiểm tra thứ ba về trinh sát và ẩn nấp trong bóng tối do sát thủ Chloe hướng dẫn. Cô nàng Cá voi dạy họ cách nín thở, điều hòa nhịp tim và hòa mình vào bóng đêm. Thế nhưng, trong trò chơi trốn tìm sinh tử này, Chloe hoàn toàn bị khuất phục trước đôi tai chó nhạy bén trên đỉnh đầu của ba mẹ con. Dù nàng cá voi sát thủ có ẩn nấp kỹ đến đâu dưới gầm bàn hay sau các kiện hàng, chiếc mũi thính nhạy của Tô và đôi tai cún của Mi vẫn phát hiện ra ngay lập tức. Sau mỗi buổi tập, Chloe lại thích thú ôm chầm lấy Ki và Tô, dụi má vào đôi tai mềm mại của họ mà khen ngợi không ngớt.

Bài kiểm tra thứ tư về lễ nghi và lòng phục tùng do chính Tổng tư lệnh La+ chủ trì. Cô bắt ba mẹ con học thuộc lòng các danh hiệu hoa mỹ của mình, đồng thời huấn luyện cách bóc bánh snack khoai tây sao cho không bị vỡ vụn và cách rót nước ngọt có ga không bị trào bọt. Mi với sự nhẫn nại của một người mẹ luôn hoàn thành mọi yêu cầu một cách chu đáo, khiến cô nàng sừng quỷ nở mày nở mặt vô cùng.

Thế nhưng, người để lại ấn tượng sâu đậm và mang đến sự ấm áp nhiều nhất cho ba mẹ con lại chính là thành viên thứ năm: Lui.

Là người phụ trách hậu cần và điều hành thực tế của holoX, cô là người có tính cách dịu dàng, chu đáo và tốt bụng nhất trong số năm người. Cô luôn là người chuẩn bị từng bữa ăn nóng hổi hợp khẩu vị cho ba mẹ con, tỉ mỉ may vá lại những chỗ sứt chỉ trên trang phục huấn luyện, và ân cần băng bó từng vết trầy xước nhỏ trên tay Tô sau các buổi tập kiếm.

Một buổi tối muộn cuối tháng Chín, khi căn cứ ngầm đã chìm vào yên ắng, Lui mang một khay trà hoa cúc nóng và đĩa bánh quy bơ vào phòng nghỉ của ba mẹ con.

Nhìn thấy Mi đang ngồi bên khung cửa sổ thông gió nhìn xa xăm về phía nam thành phố - nơi ánh sáng mờ ảo của Thị trấn Điện tử phản chiếu lên những đám mây đêm - Lui khẽ đặt tách trà xuống bàn, dịu dàng ngồi xuống bên cạnh: ``Uống chút trà nóng cho ấm người đi, Mi. Dạo này gió mùa bắt đầu về rồi, đêm ở tầng hầm lạnh lắm đấy.''

Cô Chó mẹ khẽ giật đôi tai cún trên đầu, quay lại mỉm cười đón lấy tách trà bằng hai tay: ``Cảm ơn cô, Lui-sama. Cô lúc nào cũng chu đáo với mẹ con tôi như vậy.''

Lui mỉm cười, đưa tay vén lọn tóc vướng bên má. Cô nhìn ba mẹ con đang quây quần bên nhau, ánh mắt thoáng chút tò mò nhưng đầy vẻ thấu cảm: ``Này Mi\dots\ Tôi luôn thắc mắc, trước khi đến đây, cuộc sống của ba mẹ con ở Kawachi như thế nào vậy? Những con người ở căn xưởng cơ khí ngoại ô ấy\dots\ họ đối xử với mọi người ra sao?''

Nhắc đến gia đình Hikawa, đôi mắt của Mi bỗng ánh lên một luồng sáng ấm áp kỳ lạ. Giọng cô chùng xuống, chậm rãi và đầy trìu mến: ``Họ là những người nhân hậu nhất trần đời, Lui-sama. Ryujin-sama dạn dày sương gió, cả đời lấm lem dầu máy nhưng sớm nào cũng chia cho tôi nửa ổ bánh mì nóng. Cậu chủ - Kuon - thì chăm sóc tôi từ năm cậu ấy mới mười hai tuổi, những năm tháng gia đình chia cắt đầy cô quạnh, hai bố con và tôi nương tựa vào nhau mà sống\dots\ Rồi còn có Ojou (Kaiho), Suzuno-sama, Ryuuhou-sama và Naomi-sama\dots\ Họ không bao giờ xem chúng tôi là loài vật nuôi thấp kém. Với họ, chúng tôi là người thân trong gia đình.''

Nàng Chim ưng lặng im lắng nghe, đôi mắt màu xanh nước biển ánh lên vẻ xúc động sâu sắc. Để hiểu thêm về những con người đặc biệt ấy, cô nhẹ nhàng hỏi han thêm: ``Nghe ấm áp thật đấy\dots\ Thế họ sinh vào ngày nào? Trong một gia đình đông đúc như vậy, chắc ngày sinh nhật của mỗi người rộn ràng lắm nhỉ?''

Mi chưa kịp trả lời thì cô con gái út Ki đang ngồi bên cạnh đã khẽ mở đôi mắt trong veo, cất giọng đọc vanh vách bằng trí nhớ siêu phàm:

``Ryujin-sama sinh ngày hai mươi mốt tháng Mười Một. Kaien-sama sinh ngày hai mươi sáu tháng Tư. Suzuno-sama và Ryuuhou-sama sinh cùng ngày bốn tháng Tư. Bé út Naomi sinh ngày hai mươi chín tháng Năm.''

Ki dừng lại một nhịp, đôi tai cún khẽ rung rinh: ``Và hai người con lớn, cặp song sinh Kuon và Kaiho\dots\ sinh vào ngày mười bảy tháng Bảy.''

Lui chớp mắt, ngạc nhiên trước trí nhớ tuyệt đỉnh của cô bé, rồi khẽ lẩm bẩm ghi nhớ vào trong đầu: ``Ngày mười bảy tháng Bảy à\dots? Vậy là cũng sắp tới rồi nhỉ.''

Cô mỉm cười, vỗ nhẹ lên mu bàn tay của Mi trước khi đứng dậy rời phòng. Nàng phó tướng hoàn toàn không thể ngờ rằng, những mẩu thông tin tưởng chừng như một cuộc trò chuyện tâm tình đêm muộn ấy, sau này lại trở thành căn cứ để cô vô tình buột miệng nói ra trước mặt Kuon tại văn phòng \hll\ trong một ngày sinh nhật đầy bất ngờ.

Sau khi cánh cửa phòng khép lại, căn phòng chỉ còn lại ba mẹ con dưới ánh đèn ngủ vàng nhạt.

Tô ngả đầu vào lòng mẹ, chiếc đuôi dài khẽ ngoe nguẩy: ``Mẹ ơi\dots\ bao giờ chúng ta mới được về nhà? Con nhớ khoảng sân xưởng, nhớ mùi cơm trộn thịt băm của Suzuno-sama quá\dots''

Mi khẽ siết chặt bờ vai của hai đứa con, giọng thì thầm nhưng vô cùng kiên định: ``Sắp rồi các con ạ. Chúng ta mang ơn cứu mạng của họ, nhưng mái ấm thực sự của chúng ta không thuộc về bóng tối của tầng hầm này. Hãy kiên nhẫn thêm một chút nữa. Cơ hội sắp đến rồi.''

Suốt hơn một tháng qua, ngoài mặt ba mẹ con luôn tỏ ra vâng lời, hoàn thành trọn vẹn mọi mệnh lệnh của holoX. Nhưng trong sâu thẳm tâm can, lòng trung thành nguyên bản của họ chưa từng một giây lay chuyển.

Họ chỉ đang nhẫn nại chờ đợi một thời cơ duy nhất để phá vỡ xiềng xích.

Và thời cơ ấy cuối cùng đã xuất hiện.

\ct

Một ngày giữa tháng Tám.

Để kiểm tra toàn diện khả năng tác chiến phối hợp của bộ ba cận vệ trước khi chính thức triển khai kế hoạch bành trướng thế giới, Tổng tư lệnh La+ Darknesss quyết định tổ chức một bài sát hạch thực địa quy mô lớn.

Địa điểm được lựa chọn là một vùng đất vô cùng đặc thù nằm ở rìa ngoài hẻo lánh của Thị trấn Điện tử: một bãi cát hoang vu trải dài tiếp giáp giữa sa mạc cằn cỗi và bờ biển lộng gió. Không gian nơi đây mênh mông ngút ngàn, bốn bề chỉ có cát trắng, sóng biển cuộn trào và những mỏm đá nham nhở - hoàn toàn cách biệt với khu dân cư. Đây cũng chính là bãi cát sa mạc ven biển mà nàng Iroha sau này sẽ có cuộc so tài kiếm đạo chấn động cùng Quỷ sừng Ayame.

Gió biển thổi ràn rạt, cuốn theo những hạt cát mịn bay là là trên mặt đất.

Trên một đài quan sát dã chiến dựng bằng giàn giáo thép kiên cố ở rìa cồn cát cao, La+ đứng chống nạnh, tay cầm chiếc ống dòm quân sự to tướng, bên cạnh là Takane Lui đang theo dõi các thông số truyền về từ thiết bị bay không người lái.

Nàng Tổng tư lệnh vung chiếc quạt giấy, cất giọng sang sảng qua micro phát thanh dã chiến: ``Hỡi các chiến binh của holoX! Bài kiểm tra thực chiến cuối cùng chính thức bắt đầu! Phe tấn công: Bộ ba nhà Chó do Mi chỉ huy! Phe phòng thủ: Liên quân tinh nhuệ gồm Kazama Iroha, Sakamata Chloe và bẫy công nghệ của Hakui Koyori! Mục tiêu của phe tấn công: Vượt qua toàn bộ phòng tuyến cát và chạm vào lá cờ chỉ huy trên đài quan sát trong vòng ba mươi phút! Bắt đầu!''

*TÙNG! TÙNG! TÙNG!*

Tiếng trống lệnh dã chiến vang lên giòn giã.

Ở phía bên kia bãi cát, Iroha đứng thủ thế ở trung lộ, tay cầm thanh kiếm gỗ dài, ánh mắt sắc lẹm như một võ sĩ thực thụ. Chloe nấp sau một mỏm đá ngầm ven biển, toàn thân ẩn mình trong lớp áo choàng ngụy trang cát. Còn Koyori thì nấp sau một gờ đá cao, điều khiển hàng loạt bẫy cát lún và máy phóng lưới điện sinh học.

Ba mẹ con đứng thành hình tam giác giữa bãi cát trắng mênh mông.

Đôi tai cún trên đỉnh đầu của Ki khẽ xoay tròn, đón nhận từng luồng gió biển và âm thanh dao động của máy móc. Đôi mắt vàng nâu to tròn phẳng lặng của cô bé lướt nhanh qua địa hình, bộ não siêu máy tính lập tức phân tích hàng triệu dữ liệu trong nháy mắt.

Ki khẽ ghé sát tai mẹ và chị gái, thì thầm bằng mật ngữ cực nhỏ: ``Gió nam cấp bốn, tầm nhìn bị hạn chế bởi cát bay. Góc quét radar của Koyori có điểm mù ba mươi độ ở hướng chín giờ ven mép sóng. Phòng tuyến của Iroha rất vững nhưng chân trụ trên cát lún kém linh hoạt. Đằng sau đài quan sát của La+ năm trăm mét là đường lộ liên tỉnh dẫn ra mạng lưới giao thông của thành phố.''

Mi nhìn con gái út, ánh mắt người mẹ lóe lên một tia sáng quyết đoán: ``Kế hoạch thế nào, Ki?''

``Chúng ta sẽ dùng trận giả này\dots\ để biến thành một cuộc đào thoát thực sự.''

Ánh mắt của Tô sáng rực lên. Chiếc đuôi dài phía sau quẫy mạnh trong sự phấn khích tột cùng: ``Tô hiểu rồi! Cứ để phần càn quét cho chị!''

``Lên thôi, các con!''

Mi phát lệnh. Ba bóng dáng mặc đồng phục xám tro lập tức phóng vụt đi trên mặt cát với tốc độ xé gió.

Koyori trên gờ đá giật mình kích hoạt hệ thống bẫy: ``Bẫy lưới điện sinh học số một, số hai! Kích hoạt!''

Hàng chục quả lựu đạn phóng lưới bắn ra từ lòng đất, nhắm thẳng vào hướng di chuyển của ba mẹ con. Thế nhưng, Ki đã tính toán chuẩn xác quỹ đạo từ trước. Cô bé khẽ hô: ``Trái mười bước! Nhảy!''

Cả ba người đồng loạt lạng người sang trái, tránh trọn vẹn những tấm lưới điện trong gang tấc.

Ngay sau đó, Tô lấy đà dậm mạnh chân xuống nền cát. Với sức mạnh cơ bắp gấp hai mươi lần người thường, cú dậm chân của cô nàng hất tung một bức tường cát khổng lồ cao tới bốn mét, tạt thẳng vào vị trí mai phục của Chloe ven bờ biển.

Nàng Cá voi sát thủ bị bức tường cát ập xuống bất ngờ, chiếc áo choàng ngụy trang bị cuốn bay, cô nàng ho sặc sụa: ``Khụ\dots\ khụ! Cát vào hết mồm rồi! Đâu ra mà khỏe dữ vậy trời!''

Tô không dừng lại. Cô nàng lao thẳng về phía trung lộ nơi Iroha đang trấn giữ.

Iroha quát lớn, vung kiếm gỗ tung ra một đường chém hoàn mỹ theo kiếm phái Kazama: ``\textbf{Nhận lấy này, de gozaru!}''

Thế nhưng, Tô không thèm né tránh. Cô nàng đưa cả hai cánh tay bọc giáp bảo hộ bắt trọn lấy thân kiếm của Iroha, rồi mượn đà xoay người một vòng tròn đẹp mắt, nhấc bổng cả người lẫn kiếm của nàng samurai ném văng vào một đụn cát mềm cách đó mười mét.

Iroha rơi lún sâu vào đụn cát, hai chân chổng lên trời, thanh kiếm văng ra xa: ``Oái! Sức mạnh quái vật gì thế này de gozaruuu!''

Chỉ trong vòng chưa đầy hai phút, toàn bộ hệ thống phòng thủ của holoX đã bị bẻ gãy hoàn toàn.

Trên đài quan sát, La+ Darknesss nhảy cẫng lên vỗ tay tán thưởng: ``Hay lắm! Tuyệt đỉnh! Mau chạm vào lá cờ chỉ huy đi!''

Thế nhưng, điều kỳ lạ đã xảy ra.

Thay vì lao lên bậc thang của đài quan sát để cướp cờ chiến thắng, ba mẹ con bỗng đổi hướng đột ngột.

Ki rút từ thắt lưng ra hai quả lựu đạn khói màu hồng mà Koyori từng đưa cho họ làm đạo cụ diễn tập, ném mạnh xuống chân đài quan sát.

*BÙM!!*

Một đám khói hồng khổng lồ bốc lên mù mịt, bao trùm toàn bộ đài quan sát và che khuất hoàn toàn tầm nhìn của máy bay không người lái.

``Oái! Khói ở đâu ra thế này?!'' La+ sặc sụa ho, hai tay xua xua làn khói dày đặc.

Lui giật mình nhìn vào màn hình cảm biến nhiệt: ``Khoan đã! Tín hiệu của ba người họ\dots\ họ không tiến về đài quan sát! Họ đang lao thẳng ra bờ biển hướng về đại lộ cao tốc!''

Giữa làn khói bụi mịt mù, Tô nhanh như chớp cõng mẹ Mi lên lưng. Mi vòng hai tay ôm chặt lấy cổ con gái, còn bàn tay phải của Tô thì nắm chặt lấy cổ tay nhỏ nhắn của Ki.

Chiếc đuôi dài của Tô quẫy mạnh lấy thăng bằng, đôi tai chó tam giác trên đỉnh đầu dán chặt vào sau gáy để giảm tối đa lực cản không khí. Với thể lực dũng mãnh vô song, cô nàng phóng như bay qua những cồn cát nhấp nhô, mỗi bước chạy dài tới bốn mét, để lại sau lưng những vệt cát cuộn trào như một cơn lốc xoáy.

Ki vừa chạy theo nhịp bước của chị vừa liên tục dùng thính giác và tư duy tính toán dẫn đường: ``Thẳng phía trước ba trăm mét! Vượt qua rặng phi lao là con đường dẫn ra quốc lộ số năm! Không có người truy đuổi!''

Khi làn khói hồng tan biến dần dưới ngọn gió biển, trên bãi cát sa mạc hoang vu chỉ còn lại những dấu chân dài xa tít tắp hướng về phía chân trời.

Koyori từ trên gờ đá trèo xuống, hai tai sói rũ cụp, ôm chiếc máy tính bảng gào lên thảm thiết: ``Họ\dots\ họ chạy trốn thật rồi! Mẫu vật vô giá của tôi!''

Iroha lóp ngóp bò ra khỏi đụn cát, vừa phủi cát trên đầu vừa ngơ ngác: ``Chuyện\dots\ chuyện gì vừa xảy ra thế de gozaru?''

La+ ôm đầu hét toáng lên, hai chiếc sừng quỷ rung bần bật vì uất ức: ``Trời đất ơi! Ba nữ cận vệ của ta! Đạo quân chinh phục thế giới của ta! Sao chúng nó lại dám bỏ trốn chứ?! Đuổi theo mau!''

Lui nhìn vào màn hình cảm biến trên tay, khẽ lắc đầu: ``Tốc độ di chuyển của ba người họ đã vượt quá tám mươi cây số một giờ. Hướng di chuyển dẫn thẳng vào khu trung tâm đô thị phức tạp, rất khó để triển khai vũ khí hạng nặng mà không bị chú ý.''

La+ giậm chân bình bịch trên đài quan sát, chiếc quạt giấy đập liên hồi vào thành sắt: ``Mặc kệ! Lên xe truy đuổi ngay lập tức! Phải bắt bằng được miếng mồi ngon đó về đây cho ta!''

Thế là cả năm thành viên holoX cuống cuồng nhảy lên xe dã chiến, bắt đầu một cuộc truy lùng ráo riết xuyên đêm trên khắp các ngả đường của Thị trấn Điện tử.

\ct

Băng qua ranh giới bãi cát sa mạc ven biển, ba mẹ con lao mình vào màn đêm vô định của vùng ngoại vi Thị trấn Điện tử.

Ở thế giới bên này, các khu phố công nghệ ngập tràn những bảng hiệu neon đa sắc, những tòa cao ốc kính chọc trời và mạng lưới đường xá chằng chịt xa lạ. Không có ngọn cỏ lau quen thuộc của đồng quê Kawachi, cũng không có mùi đất ẩm nồng sau cơn mưa rào của quê nhà Etsunan. Ba mẹ con dừng chân dưới bóng râm của một con hẻm hẹp, đưa mắt nhìn quanh giữa một vùng không gian hoàn toàn xa lạ.

Đúng lúc đang mất phương hướng giữa muôn trùng ngã rẽ, Tô bỗng giật nảy đôi tai chó tam giác trên đỉnh đầu.

Chiếc mũi thon nhỏ của cô nàng khẽ hếch lên vài nhịp trong làn gió đêm thoang thoảng: ``Khịt\dots\ khịt khịt\dots''

Một mùi hương kỳ lạ bỗng xộc thẳng vào khứu giác nhạy bén của Tô. Đó không phải là mùi dầu mỡ công nghiệp của Thị trấn Điện tử, mà là mùi nước xả vải dịu nhẹ quen thuộc quyện cùng hương thơm của một người phụ nữ nghiêm nghị, quyết đoán.

Tô tròn xoe mắt, nắm lấy tay áo Mi: ``Mẹ ơi\dots\ Mùi này\dots\ là mùi của Kaiho-sama!''

Mi và Ki cùng giật mình nhìn theo hướng chỉ tay của Tô.

Cách con hẻm chừng năm mươi mét, nép mình bên hông một tòa nhà phức hợp là một cánh cửa ngách hẹp bằng kim loại sơn xám. Cánh cửa khép hờ, bên trong hắt ra ánh đèn vàng ấm áp chiếu rọi lên không gian mở hai tầng rộng rãi. Qua khe cửa kính mờ, có thể nhìn thấy những bức tường sơn màu tím nhạt nhám thô ráp, những cột thạch cao trắng tinh khôi, sàn lát gỗ lớn màu nâu sẫm và một quầy bar gỗ tự nhiên với hàng ghế đẩu chân cao bọc nệm đỏ. Trên chiếc bàn ăn dài bằng gỗ thô ở trung tâm căn phòng, một vài tập tài liệu quản lý còn lưu lại những nét chữ viết tay dứt khoát, sắc sảo.

Đó chính là căn mancave - trụ sở sinh hoạt chung của các chàng trai nhóm \hls, nơi mà Kaiho đã đến làm Tổng quản lý để thiết lập kỷ luật thép và giám sát họ.

Dù lúc này căn mancave đã vắng lặng sau giờ giới nghiêm, nhưng mùi hương của Kaiho vẫn còn vương vấn đậm đặc quanh lối đi phụ.

Lần theo vệt mùi hương quen thuộc ấy, ba mẹ con rảo bước qua những con ngõ nhỏ của Thị trấn Điện tử. Vệt mùi của Kaiho dẫn họ đi ngược về một tòa nhà văn phòng cao chừng chín, mười tầng nằm khiêm tốn giữa những khối nhà chọc trời. Trên mặt tiền của tòa nhà, dòng chữ logo \textit{hololive} màu xanh lam phát sáng nổi bật trong đêm.

Đứng trước sảnh chính của tòa nhà văn phòng COVER Corp, người mẹ bỗng khựng lại.

Chiếc mũi của Mi khẽ rung rinh. Đôi tai cún phủ lông tơ tro bạc trên đỉnh đầu cô dựng đứng lên. Giữa muôn vàn mùi hương hỗn tạp của thành phố Điện tử, một luồng hương vị mộc mạc, trầm ấm đã ăn sâu vào từng tế bào tủy sống của Mi suốt mười bốn năm qua bỗng chốc ùa về, rõ rệt và thiêng liêng đến nghẹn ngào.

Đó là mùi cơ thể ấm áp của người cậu bé mười hai tuổi năm nào, mùi chiếc áo thun đỏ xuề xòa và mùi mồ hôi sau những giờ làm việc mệt nhoài.

Mùi của cậu chủ Kuon.

``Là Cậu chủ\dots\ Cậu chủ từng ở đây!''

Đôi mắt Mi đỏ hoe, ngấn lệ. Vệt mùi hương của Kuon lưu lại trên bậc thềm sảnh dẫn thẳng ra một đại lộ lớn, hướng về phía nhà ga trung tâm của khu vực Shinjuku.

Ba mẹ con vội vã bám theo vệt hương ấy, chạy một mạch tới cửa ngõ nhà ga tàu điện ngầm khổng lồ.

Lúc này đã quá nửa đêm, các chuyến tàu thường nhật đã ngừng hoạt động, các cửa xoay soát vé đóng kín mít. Nhưng ở sâu bên dưới tầng hầm sân ga, một cánh cổng vòm ánh sáng mờ ảo mang tên ``Cổng Warp liên quốc gia'' vẫn đang duy trì luồng từ trường không gian xanh ngắt dành riêng cho các tuyến đường sắt liên chiều.

Ki bước lại gần bảng sơ đồ tuyến đường sắt ngầm gắn bên vách kính. Đôi mắt vàng phẳng lặng của cô bé lướt qua hàng loạt nút giao thông phức tạp. Bộ não siêu tính toán lập tức kích hoạt, xử lý toàn bộ các phương trình trắc địa và bản đồ địa lý liên chiều chỉ trong vài giây.

Cô quay lại nhìn mẹ và chị gái, đôi tai cún khẽ giật nhẹ, cất giọng rành rọt: ``Tuyến đường sắt ngầm này chạy xuyên qua cổng warp không gian, kết nối trực tiếp Thị trấn Điện tử này với thành phố Kawachi của đất nước Etsunan. Điểm dừng đầu tiên sau cánh cổng là Ga \ruby{Kinhasu}{Kim Liên}, nằm ngay sát cổng trường đại học KUCE mà Cậu chủ và Cô chủ theo học.''

Tô nghe vậy liền sốt sắng: ``Vậy chúng ta đến đó rồi trèo lên mặt đất luôn sao em?''

Ki lắc đầu, ngón tay nhỏ nhắn chỉ tiếp theo vệt đường ray trên sơ đồ: ``Không nên, chị Tô. Từ Ga Kinhasu về tới nhà mình ở phố Anwa cách nhau rất nhiều cây số, nếu trèo lên mặt đất đi bộ giữa đêm sẽ dễ bị tuần tra phát hiện hình dạng tai và đuôi cún. Theo tính toán của em, tuyến ngầm này chạy thẳng tắp dưới lòng đất qua khu trung tâm, điểm kế tiếp thuận lợi nhất chính là Ga Vành đai 2. Nhà ga đó nằm ngay nút giao cầu lớn, trồi lên sẽ là ngã tư đường. Đi men theo đường dọc bờ sông là có thể đi thẳng vào lòng phố Anwa của phường mình.''

Kế hoạch hoàn hảo và chuẩn xác đến từng chi tiết khiến Mi và Tô đều khẽ gật đầu đồng thuận.

Không chần chừ, ba bóng hình mặc áo choàng xám tro khéo léo lách qua khe rào chắn của sân ga vắng lặng, nhảy xuống lòng đường ray tối om và tiến thẳng vào luồng ánh sáng của cánh cổng warp liên không gian.

\ct

Một luồng sáng xanh biếc lướt qua, không khí xung quanh bỗng chốc biến chuyển hoàn toàn.

Mùi kim loại và điện từ khô khốc của Thị trấn Điện tử tan biến, nhường chỗ cho luồng gió ngầm mát lành mang theo độ ẩm quen thuộc của đất trời quê hương.

Họ đã vượt qua cổng, đặt chân vào lòng đường hầm ngầm thuộc mạng lưới đường sắt thành phố Kawachi. Lúc này đường ray hoàn toàn tĩnh mịch vì các đoàn tàu đã về ga tập kết sau chuyến cuối trong ngày.

Nhờ đôi mắt sáng rõ trong bóng tối và đôi tai cún thính nhạy có thể nghe thấy từng xung động của luồng gió trong các giếng thông gió, ba mẹ con rảo bước thoăn thoắt dọc theo hai thanh ray kim loại. Chẳng mấy chốc, ánh đèn vàng mờ của sân ga ngầm Kinhasu lướt qua bên vách hầm. Đôi mũi của họ thoáng ngửi thấy mùi hương quen thuộc của sách vở và giảng đường KUCE nơi Kuon cùng Kaiho gắn bó mỗi ngày, mang lại một nguồn động lực tinh thần to lớn.

Không dừng lại ở Kinhasu, ba mẹ con tiếp tục men theo lòng hầm ngầm hun hút chạy thẳng về phía tây. Dưới sự dẫn đường chuẩn xác của Ki, họ vượt qua hàng nghìn mét đường hầm ngầm trong yên lặng, tới điểm dừng Ga Vành đai 2.

Tại chân Ga Vành đai 2, Ki dẫn mẹ và chị lách qua cửa thoát hiểm kỹ thuật, bước lên những bậc thang xoắn ốc bằng thép dẫn thẳng lên mặt đất.

Đẩy nhẹ cánh cửa thông gió nép bên lùm cây xanh, ba mẹ con đặt chân lên mặt đường nhựa phẳng lặng. Những cơn gió về đêm se lạnh mơn man da thịt, đưa hương thơm của dòng sông và cây cỏ tràn vào lồng ngực.

Họ đang đứng ngay tại ngã tư đường, nơi đầu cầu vươn qua mặt nước sông phản chiếu ánh đèn đường lấp lánh.

Từ ngã tư này, con đường Nguyễn Khang trải dài tĩnh lặng dọc theo bờ sông, vắng bóng xe cộ qua lại, phẳng phiu dẫn thẳng một mạch vào địa phận phường Yên Hòa, nơi có con phố Anwa thân thương.

Không còn xe cộ, cũng không có phương tiện hỗ trợ nào giữa đêm khuya thanh vắng, ba mẹ con quyết định cùng nhau cất bước hành bộ.

Một cuộc cuốc bộ ròng rã kéo dài suốt ba tiếng đồng hồ từ lúc rời Thị trấn Điện tử.

Đôi bàn chân trần con người mới mẻ chưa quen với việc đi bộ đường dài dần trở nên rã rời, mỏi nhừ và phồng rộp trên mặt đường nhựa lởm chởm đá dăm. Những giọt mồ hôi lăn dài trên gò má, nhưng không một ai cất lời than vãn.

Tô bước đi hiên ngang phía trước, thỉnh thoảng lại đưa cánh tay săn chắc đỡ lấy bước chân của mẹ. Ki lặng lẽ bước bên cạnh, đôi mắt luôn quan sát xung quanh để dẫn đường tránh những đoạn phố có ánh đèn pha xe tải chạy đêm. Còn Mi thì siết chặt lấy bàn tay của hai đứa con gái, trong lòng ngập tràn một niềm xúc động vô bờ bến.

Mỗi bước chân đặt xuống là một bước tiến gần hơn tới mái ấm gia đình.

\ct

Ba giờ sáng.

Bầu trời Kawachi lấp lánh muôn ngàn ánh sao thu trong trẻo. Những ngọn gió đêm mát lành thổi qua những rặng cây ven đường xao xác.

Sau ba tiếng đồng hồ hành bộ miệt mài, ba bóng người khoác áo choàng xám tro cuối cùng đã dừng lại trước ngã tư quen thuộc dẫn vào con phố Anwa.

Tô dừng bước, chiếc mũi xinh xắn khẽ hếch lên vài nhịp hít hà trong không khí. Đôi tai tam giác trên đỉnh đầu vểnh cao hết cỡ, chiếc đuôi dài sau lưng quẫy mạnh trong niềm vui sướng khôn tả: ``Mùi dầu máy quen thuộc từ xưởng cơ khí của Ryujin-sama\dots\ Mùi nước hoa cúc của Suzuno-sama từ phía căn nhà tám tầng kia rồi\dots\ Mẹ ơi! Em Ki ơi! Đúng là nhà mình đây rồi! Chúng ta đã về tới nơi rồi!''

Ki mỉm cười dịu dàng, đôi mắt vàng to tròn phản chiếu ánh đèn đường lung linh, hai bàn tay nhỏ bé ôm chặt lấy vạt áo mẹ.

Mi đứng giữa hai đứa con gái, ngẩng đầu nhìn về phía căn nhà sáu tầng quen thuộc đang chìm trong giấc ngủ bình yên nơi cuối con phố Anwa. Đôi mắt người mẹ ngấn lệ, bờ môi khẽ run run nhưng ánh lên niềm hạnh phúc thiêng liêng tột cùng:

``Ừ\dots\ Ba mẹ con mình về nhà rồi các con ạ.''

Dưới ánh trăng thu bàng bạc, ba bóng hình rảo bước nhanh hơn về phía cánh cổng sắt của khu ký túc xá, mở màn cho cuộc tao ngộ kỳ tích và chấn động sắp sửa diễn ra tại gia đình Hikawa\dots

\ct

Cùng lúc đó, tại căn cứ ngầm của Tổ chức Bí mật holoX ở Thị trấn Điện tử.

Bầu không khí trong phòng chỉ huy lúc này u ám và hỗn loạn như vừa trải qua một trận cuồng phong. Chiến dịch truy lùng ráo riết suốt đêm của cả năm thành viên hoàn toàn thất bại. Mọi manh mối về ba mẹ con đều biến mất không dấu vết ngay tại khu vực ga ngầm trung tâm.

La+ nằm vật ra chiếc ghế bành chỉ huy bọc nhung đen, hai chân đạp loạn xạ lên bàn, chiếc quạt giấy bị vò nát quăng sang một bên. Nàng Tổng tư lệnh phụng phịu gào khóc, hai chiếc sừng quỷ lắc lư bần bật vì uất ức: ``Oaaaa! Miếng mồi ngon béo bở đến thế mà lại để tuột mất! Một đứa khỏe như quái vật đấm bay cả samurai, một đứa thông minh hơn siêu máy tính lượng tử, một con chó mẹ đầy uy quyền chỉ huy! Bao nhiêu công sức nuôi nấng, dạy dỗ cả tháng trời của ta tan thành mây khói hết rồi! Giấc mộng chinh phục thế giới của ta biết tính làm sao đây?!''

Bên cạnh bàn điều khiển, Koyori ngồi sụp xuống sàn, hai tai sói hồng rũ rượi, đôi mắt ầng ậng nước nhìn chằm chằm vào dòng chữ ``SIGNAL LOST'' nhấp nháy đỏ lòm trên màn hình: ``Dữ liệu tiến hóa hoàn hảo gấp mười lần, hai mươi lần của tôi\dots\ Huyết thanh tinh chỉnh còn chưa kịp trích xuất trọn vẹn\dots\ Công trình nghiên cứu sinh học đỉnh cao nhất cuộc đời tôi bay mất rồi! Không biết sau này làm sao tìm lại được một mẫu vật kỳ diệu như thế nữa\dots''

Chloe ngồi bệt dựa lưng vào chân bàn, vừa nhồm nhoàm gặm dở một mẩu bánh mì khô khốc vừa đưa tay gãi mớ tóc bù xù còn vương cát biển. Nàng Cá voi sát thủ tặc lưỡi, cất giọng tỉnh bơ: ``Thôi bớt than vãn đi hai bà tướng. Con bé Tô đó lúc ở bãi cát chỉ dậm nhẹ một cái là bức tường cát bốn mét hất văng em rồi. Đuổi kịp nó rồi ai bắt? Để nó túm cổ quăng vèo lên trời nữa à? Tôi cuốc bộ đuổi theo cả đêm rã hết cả chân tay rồi, giờ chỉ muốn đi ngủ thôi\dots\ À mà cát dính đầy người thế này chắc phải tắm, nhưng lười quá, để mai tính.''

Iroha ngồi trên chiếc ghế gỗ dài gần cửa, tay vẫn đang nhẹ nhàng xoa nắn bả vai còn ê ẩm sau cú va chạm vào đụn cát. Nàng Samurai cẩn thận tra thanh kiếm gỗ vào bao, khẽ thở dài đầy thán phục: ``Tại hạ tu luyện kiếm đạo khổ hạnh bao năm trên núi tuyết, vậy mà cô nương ấy dùng tay không bắt trọn lưỡi kiếm rồi nhấc bổng tại hạ nhẹ như không de gozaru\dots\ Thể lực vô song, phản xạ xuất thần, lại có tiểu muội bên cạnh tính toán mọi đường đi nước bước. Với bản lĩnh ấy, nếu họ đã quyết tâm dứt áo ra đi thì dù có bắt lại được, sớm muộn căn cứ của chúng ta cũng bị họ san phẳng thôi de gozaru.''

Giữa những tiếng cằn nhằn và than thở dở khóc dở cười của bốn người đồng đội, Lui chỉ lặng lẽ đứng bên khung cửa sổ kính lớn nhìn ra bầu trời đêm của Thị trấn Điện tử.

Nàng Chim ưng khẽ nâng tách cà phê nóng, nhấp một ngụm nhỏ. Trong làn khói mờ ảo bốc lên từ chiếc tách sứ, những ký ức của hơn một tháng qua chầm chậm lướt qua tâm trí cô: ánh mắt kiên định của người mẹ Mi mỗi khi che chở cho hai đứa con, niềm vui sướng ngây thơ của Tô khi nhắc tới góc sân xưởng cơ khí cùng đĩa cơm trộn thịt băm, và giọng nói rụt rè, trân trọng của Ki khi ghi nhớ từng ngày sinh nhật của những người trong gia đình Hikawa\dots

Lui hiểu rõ hơn bất kỳ ai. Ba mẹ con không phải là những cỗ máy chiến tranh hay những công cụ phục vụ dã tâm bành trướng. Họ là những sinh mệnh có linh hồn, có trái tim, và mang một lòng trung thành nguyên bản tuyệt đối với mái ấm duy nhất của mình. Một khi sợi dây gắn kết gia đình thiêng liêng ấy đã lên tiếng dẫn lối, thì không một thế lực hay xiềng xích nào có thể cầm chân họ.

Một nụ cười bao dung, ấm áp chầm chậm nở trên khóe môi nàng phó tướng.

Lui ngước nhìn vầng trăng khuyết đang lặn dần về phương tây xa xăm, nơi có thành phố Kawachi và con phố Anwa bình yên, khẽ nâng nhẹ tách cà phê trong tay như một lời từ biệt chân thành:

``Sống cho thật khỏe mạnh nhé, ba mẹ con.''

\end{document}